\section*{Chapitre \ref{ch:b1:optical-mineralogy} --- Minéralogie optique}

\begin{solution}{exo:b1:optical-mineralogy:1}
Quartz~: $\qty{30000}{nm}\times0.009 = \qty{270}{nm}$, un blanc du premier \omterm{def:b1:optical-mineralogy:interference-colour}{ordre}.
Calcite~: $\qty{30000}{nm}\times0.172 = \qty{5160}{nm}$, bien au-delà de l'abaque~:
un blanc pâle d'\omterm{def:b1:optical-mineralogy:interference-colour}{ordre} élevé, faiblement rosé ou verdâtre.
\end{solution}

\begin{solution}{exo:b1:optical-mineralogy:2}
$I/I_{\max} = \sin^2 2\theta$~: $0$ à \ang{0}~; $\sin^2\ang{45} = 0.5$ à
\ang{22.5}~; $1$ à \ang{45}~; $0.5$ à \ang{67.5}~; $0$ à \ang{90}.
\end{solution}

\begin{solution}{exo:b1:optical-mineralogy:3}
Le \omterm{def:b1:rock-forming-minerals:groups}{grenat} est cubique, donc isotrope~: la lumière le quitte en vibrant comme elle
y est entrée, le long du polariseur, et l'analyseur l'arrête~; $\Delta = 0$ dans
toutes les positions. Un grain de quartz coupé perpendiculairement à son \omterm{def:b1:optical-mineralogy:optic-axis}{axe
optique} a lui aussi $n_1 = n_2$ pour la lumière qui se propage le long de cet
axe, donc $\Delta = 0$~: il reste sombre lui aussi. En \omterm{def:b1:field-lab-methods-1:microscope}{lumière polarisée non
analysée}, le \omterm{def:b1:rock-forming-minerals:groups}{grenat} a un \omterm{def:b1:optical-mineralogy:relief}{relief} très fort ($1.83$ contre $1.54$), le quartz
presque nul.
\end{solution}

\begin{solution}{exo:b1:optical-mineralogy:4}
$|n - 1.54|$~: almandin $0.29$, forstérite $0.11$, orthose $0.02$, quartz
$0.01$~; \omterm{def:b1:optical-mineralogy:relief}{relief} dans cet \omterm{def:b1:optical-mineralogy:interference-colour}{ordre}, du plus fort au plus faible. L'orthose a l'indice
le plus bas~: \omterm{def:b1:optical-mineralogy:relief}{relief} négatif, et son \omterm{def:b1:optical-mineralogy:relief}{liseré de Becke} sort du grain, vers la
résine, quand la mise au point est remontée.
\end{solution}

\begin{solution}{exo:b1:optical-mineralogy:5}
$\delta = 650/30000 = 0.022$. Un \omterm{def:b1:rock-forming-minerals:mineral}{minéral} vert \omterm{def:b1:optical-mineralogy:pleochroism}{pléochroïque} à clivages à
\ang{56} et \ang{124} et de \omterm{def:b1:optical-mineralogy:birefringence}{biréfringence} voisine de $0.020$~: la hornblende.
\end{solution}

\begin{solution}{exo:b1:optical-mineralogy:6}
$\tan\ang{30} = 0.5774$. Quartz~: $(0.013/1.54)\times0.5774 = \qty{0.00487}{rad}
= \ang{0.28}$. Forstérite~: $(0.111/1.54)\times0.5774 = \qty{0.0416}{rad} =
\ang{2.38}$. Rapport $0.111/0.013 = 8.5$. Exactement~: $\sin r = 1.54\sin\ang{30}/1.651
= 0.4664$, $r = \ang{27.80}$, soit une déviation de \ang{2.20}~: la formule du
premier \omterm{def:b1:optical-mineralogy:interference-colour}{ordre} surestime de \qty{8}{\percent} pour une différence d'indice aussi
grande.
\end{solution}

\begin{solution}{exo:b1:optical-mineralogy:7}
$270 + 540 = \qty{810}{nm}$, un vert du deuxième \omterm{def:b1:optical-mineralogy:interference-colour}{ordre}, lorsque les directions
lentes du grain et de la lame sont parallèles (les retards s'ajoutent)~; $|270 - 540| =
\qty{270}{nm}$, de nouveau un blanc du premier \omterm{def:b1:optical-mineralogy:interference-colour}{ordre}, lorsque la direction lente
du grain est parallèle à la direction rapide de la lame.
\end{solution}

\begin{solution}{exo:b1:optical-mineralogy:8}
En pleine épaisseur, $\Delta = \qty{30000}{nm}\times0.052 = \qty{1560}{nm}$. Du
bord ($\Delta = 0$) vers l'intérieur, la \omterm{def:b1:optical-mineralogy:retardation}{différence de marche} passe par 542 et
\qty{1074}{nm} mais pas par \qty{1605}{nm}~: deux franges rouge-violet.
\end{solution}

\begin{solution}{exo:b1:optical-mineralogy:9}
L'\omterm{def:b1:optical-mineralogy:extinction}{angle d'extinction} mesuré sur un grain dépend de la façon dont le grain est
coupé~: seules les sections d'une orientation particulière montrent le plus grand
angle, et les autres en montrent moins, jusqu'à zéro. La plus grande valeur,
environ \ang{37}, est celle qu'il faut comparer~: elle correspond au diopside
(environ \ang{38}), non à la hornblende (au plus \ang{19}).
\end{solution}

\begin{solution}{exo:b1:optical-mineralogy:10}
Avec l'analyseur selon $x$, les composantes en $x$ des deux ondes sont
$E_0\cos\theta\cdot\cos\theta$ et $-E_0\sin\theta\,e^{-i\varphi}\cdot(-\sin\theta)$,
donc $A = E_0(\cos^2\theta + \sin^2\theta\,e^{-i\varphi})$ et
\[
  |A|^2 = E_0^2\bigl(\cos^4\theta + \sin^4\theta + 2\sin^2\theta\cos^2\theta\cos\varphi\bigr)
  = E_0^2\bigl(1 - 2\sin^2\theta\cos^2\theta\,(1 - \cos\varphi)\bigr),
\]
en utilisant $\cos^4\theta + \sin^4\theta = 1 - 2\sin^2\theta\cos^2\theta$. Comme
$2\sin^2\theta\cos^2\theta = \frac12\sin^2 2\theta$ et $1 - \cos\varphi =
2\sin^2(\varphi/2)$, $I_\parallel = I_0(1 - \sin^2 2\theta\sin^2(\pi\Delta/\lambda))$.
À chaque longueur d'onde, $I_\parallel + I_\perp = I_0$~: les couleurs entre
polariseurs parallèles sont les complémentaires de celles entre \omterm{def:b1:field-lab-methods-1:microscope}{polariseurs
croisés}. Un grain de \omterm{def:b1:optical-mineralogy:retardation}{différence de marche} \qty{550}{nm}, magenta entre
\omterm{def:b1:field-lab-methods-1:microscope}{polariseurs croisés}, est vert jaunâtre pâle entre polariseurs parallèles.
\end{solution}

\begin{solution}{exo:b1:optical-mineralogy:11}
(a) $\Delta/\lambda$ va de $5000/700 = 7.14$ à $5000/400 = 12.5$~; il passe par
les entiers 8 à 12~: cinq zéros. (b) Les zéros sont en $\lambda_k = \Delta/k$~;
deux zéros successifs sont distants de $\Delta/k - \Delta/(k + 1) =
\Delta/(k(k + 1)) \approx \lambda^2/\Delta$, soit $550^2/5000 =
\qty{60.5}{nm}$ près de \qty{550}{nm}. (c) Cinq bandes sombres étroites réparties
sur le spectre retirent un peu de chaque couleur, et le reste passe avec un poids
moyen d'un demi~: l'œil voit une teinte pâle, presque neutre.
\end{solution}

\begin{solution}{exo:b1:optical-mineralogy:12}
Sur un tour complet, la moyenne de $\sin^2 2\theta$ est $\frac1{2\pi}\int_0^{2\pi}\sin^2 2\theta\,
\mathrm d\theta = \frac12$, donc l'intensité moyenne est
$\frac12 I_0\sin^2(\pi\Delta/\lambda)$. Elle calcule $\delta = \Delta/30000$
au lieu de $\Delta/25000$ (en nm)~: toutes ses valeurs valent $25/30 = 0.83$ des
valeurs vraies, trop basses de \qty{17}{\percent}.
\end{solution}

\begin{solution}{pb:b1:optical-mineralogy:1}
\textbf{1.} B~: les bords sombres signifient un \omterm{def:b1:optical-mineralogy:relief}{relief} fort~; les limites à peine
visibles de A signifient un \omterm{def:b1:optical-mineralogy:relief}{relief} faible.

\textbf{2.} A~: $(1.544 + 1.553)/2 = 1.5485$, $|n - n_m| = 0.0085$. B~:
$(1.664 + 1.672 + 1.694)/3 = 1.6767$, $|n - n_m| = 0.1367$.

\textbf{3.} Les déviations sont proportionnelles à $|n - n_m|$~: $0.1367/0.0085 =
16.1$.

\textbf{4.} B a un indice plus élevé que la résine~: \omterm{def:b1:optical-mineralogy:relief}{relief} positif.

\textbf{5.} \omterm{def:b1:optical-mineralogy:pleochroism}{Pléochroïsme}~; biotite.

\textbf{6.} Deux clivages à \ang{56} et \ang{124}~: une \omterm{def:b1:rock-forming-minerals:groups}{amphibole}.

\textbf{7.} $\qty{30000}{nm}\times0.009 = \qty{270}{nm}$~: blanc du premier \omterm{def:b1:optical-mineralogy:interference-colour}{ordre}.

\textbf{8.} Chaque grain est coupé selon sa propre orientation et montre une
valeur de $|n_1 - n_2|$ comprise entre 0 et $\delta$~; seuls les grains coupés
parallèlement à l'\omterm{def:b1:optical-mineralogy:optic-axis}{axe optique} montrent $\delta$ lui-même. La plus haute teinte
parmi de nombreux grains est celle à utiliser.

\textbf{9.} $d = 342/0.009 = \qty{38000}{nm} = \qty{38}{\micro m}$.

\textbf{10.} $(38 - 30)/30 = \qty{26.7}{\percent}$.

\textbf{11.} $\sin^2(\pi\times342/550) = 0.861$; at \qty{30}{\micro m},
$\sin^2(\pi\times270/550) = 0.999$.

\textbf{12.} $\qty{38000}{nm}\times0.008 = \qty{304}{nm}$ à
$\qty{38000}{nm}\times0.013 = \qty{494}{nm}$~: du blanc à l'orange-rouge, tout en
dessous de la fin du premier \omterm{def:b1:optical-mineralogy:interference-colour}{ordre} à \qty{542}{nm}.

\textbf{13.} Entre 1074 et \qty{1605}{nm}~: troisième \omterm{def:b1:optical-mineralogy:interference-colour}{ordre}.

\textbf{14.} $\delta = 1140/38000 = 0.030$.

\textbf{15.} $1140/30000 = 0.038$, ce qui aurait suggéré une \omterm{def:b1:rock-forming-minerals:groups}{olivine}. Une \omterm{def:b1:rock-forming-minerals:groups}{olivine}
magnésienne et du quartz réagissent pour donner du \omterm{def:b1:rock-forming-minerals:groups}{pyroxène},
\ce{Mg2SiO4 + SiO2 -> 2 MgSiO3}~; une roche qui a cristallisé avec du quartz ne
devrait pas garder les deux.

\textbf{16.} $0.030$ est la \omterm{def:b1:optical-mineralogy:birefringence}{biréfringence} du diopside~; son \omterm{def:b1:optical-mineralogy:relief}{relief} fort, son
absence de couleur et ses deux clivages voisins de \ang{87} concordent~: B est un
clinopyroxène, le diopside.

\textbf{17.} $\qty{30000}{nm}\times0.030 = \qty{900}{nm}$~: jaune du deuxième \omterm{def:b1:optical-mineralogy:interference-colour}{ordre}.

\textbf{18.} De \qty{8}{\micro m}~; le quartz ne montre alors au plus que
\qty{270}{nm}, blanc du premier \omterm{def:b1:optical-mineralogy:interference-colour}{ordre}.

\textbf{19.} \omterm{def:b1:optical-mineralogy:twin}{Macle polysynthétique}~: D est du \omterm{def:b1:rock-forming-minerals:groups}{plagioclase}.

\textbf{20.} La composition du \omterm{def:b1:rock-forming-minerals:groups}{plagioclase} (sa teneur en anorthite).

\textbf{21.} Une \omterm{def:b1:rock-forming-minerals:groups}{amphibole} d'\omterm{def:b1:optical-mineralogy:extinction}{angle d'extinction} \ang{17}~: la hornblende
(\ang{15}--\ang{19}), non un \omterm{def:b1:rock-forming-minerals:groups}{pyroxène} (\ang{38} pour le diopside).

\textbf{22.} $\qty{38000}{nm}\times0.020 = \qty{760}{nm}$~: un vert du deuxième
\omterm{def:b1:optical-mineralogy:interference-colour}{ordre}.

\textbf{23.} Il est coupé perpendiculairement à un \omterm{def:b1:optical-mineralogy:optic-axis}{axe optique}~: ses deux indices
sont égaux pour la lumière qui se propage le long de l'axe, $\Delta = 0$, et il
reste sombre comme un grain isotrope.

\textbf{24.} A quartz, B diopside, C biotite, D \omterm{def:b1:rock-forming-minerals:groups}{plagioclase}, E hornblende~: une
roche magmatique intermédiaire avec un peu de quartz, comme une diorite quartzique.

\textbf{25.} La section a \textbf{\qty[text-series-to-math=true]{38}{\micro m}} d'épaisseur.
\end{solution}
